% Source of input/bench/en_book/kb/electronic-computing/stored_program.pdf
% Written for the AutoGenBook benchmark. Original text, Apache 2.0 (repository licence).
% Build: lualatex stored_program.tex (twice). Running header and page-number footer are
% deliberate: the new extractor must drop them.
\documentclass[11pt,a4paper]{article}
\usepackage{fontspec}
\usepackage[margin=28mm]{geometry}
\usepackage{fancyhdr}
\usepackage{booktabs}
\usepackage{array}
\pagestyle{fancy}
\fancyhf{}
\fancyhead[L]{The Stored-Program Computer}
\fancyhead[R]{Benchmark corpus}
\fancyfoot[C]{\thepage}
\title{The Stored-Program Computer: From Theory to Transistors}
\author{AutoGenBook benchmark corpus}
\date{}
\begin{document}
\maketitle
\thispagestyle{fancy}

\section{A theory of computation}

In 1936 the British mathematician Alan Turing published a paper on computable
numbers in which he described an abstract machine that reads and writes
symbols on an unbounded tape according to a finite table of rules. Turing
showed that a single universal machine could imitate any other such machine
if the description of that machine was written on its tape. The universal
machine anticipated the idea that a program is itself a kind of data. In the
same year the American logician Alonzo Church reached equivalent conclusions
with a different formalism, and the claim that these models capture
everything that can be computed mechanically became known as the
Church--Turing thesis.

\section{Electronic machines of the war years}

The first electronic digital devices used vacuum tubes as switches, because
tubes could change state thousands of times faster than the relays of
electromechanical calculators. Between 1939 and 1942 John Atanasoff and
Clifford Berry built at Iowa State College a machine for solving systems of
linear equations. It computed in binary and stored intermediate values on
rotating drums of capacitors, but it was not programmable.

In Britain, the engineer Tommy Flowers designed Colossus to help break the
Lorenz cipher used by the German high command. The first Colossus was
operational at Bletchley Park in early 1944. It was programmed by switches
and plugs, and ten machines were in use by the end of the war. Their
existence remained secret until the 1970s.

\section{ENIAC}

ENIAC, the Electronic Numerical Integrator and Computer, was built at the
University of Pennsylvania by John Mauchly and J. Presper Eckert and was
completed in 1945. It was designed to calculate artillery firing tables for
the United States Army. The machine contained about 17,468 vacuum tubes,
weighed around 27 tonnes and occupied a large room.

ENIAC was fast, performing about five thousand additions per second, but it
had no stored program. To change the problem it solved, operators had to
reconnect cables on plugboards and set thousands of switches, a process that
could take days. The women who programmed it, among them Kathleen McNulty,
Jean Jennings and Frances Elizabeth Holberton, worked out the wiring from
the mathematical description of each problem.

\section{The EDVAC report and the stored-program concept}

Even before ENIAC was finished, its designers planned a successor that would
avoid the rewiring problem. In June 1945 the mathematician John von Neumann,
who had joined the discussions as a consultant, wrote the First Draft of a
Report on the EDVAC. The report described a machine whose instructions were
held in the same memory as the data, encoded as numbers. A program could
then be loaded as quickly as data, and the machine could even modify its own
instructions.

The organisation described in the report consisted of a central arithmetic
unit, a control unit, a memory, and input and output devices. It became
known as the von Neumann architecture, although the ideas were developed
jointly with Eckert and Mauchly, who disputed the attribution. The
separation between a processor that fetches instructions and a memory that
holds both instructions and data is still the basis of most computers.

\section{The first stored programs}

The first machine to run a program stored in its electronic memory was the
Manchester Small-Scale Experimental Machine, nicknamed the Baby, built by
Frederic Williams, Tom Kilburn and Geoff Tootill at the University of
Manchester. On 21 June 1948 it ran a program that found the highest proper
factor of a number. Its memory was a cathode-ray tube that stored bits as
charged spots on the screen, a device known as the Williams tube.

The first stored-program computer to provide a regular computing service was
EDSAC, built at the University of Cambridge under the direction of Maurice
Wilkes. It ran its first program in May 1949. The EDSAC team introduced a
library of subroutines and a simple assembler called initial orders, which
translated mnemonic instructions into machine code when a program was
loaded.

\begin{table}[h]
\centering
\begin{tabular}{>{\raggedright}p{2.8cm} >{\raggedright}p{1.6cm} >{\raggedright}p{3.6cm} >{\raggedright\arraybackslash}p{4.2cm}}
\toprule
Machine & Year & Switching technology & Programming method \\
\midrule
Colossus & 1944 & Vacuum tubes & Switches and plugs \\
ENIAC & 1945 & Vacuum tubes & Plugboards and switches \\
Manchester Baby & 1948 & Vacuum tubes & Stored program \\
EDSAC & 1949 & Vacuum tubes & Stored program with subroutines \\
\bottomrule
\end{tabular}
\caption{Early electronic computers}
\label{tab:machines}
\end{table}

\section{From vacuum tubes to transistors}

Vacuum tubes were bulky, consumed a lot of power and failed often. In
December 1947 John Bardeen and Walter Brattain, working in the group of
William Shockley at Bell Laboratories, demonstrated the point-contact
transistor, a solid-state amplifier and switch made of germanium. The three
received the Nobel Prize in Physics in 1956. By the late 1950s transistors
had replaced tubes in new computers, which became smaller, cheaper and far
more reliable.

The next step was to place many transistors on a single piece of
semiconductor. In 1958 Jack Kilby at Texas Instruments built the first
integrated circuit from germanium, and in 1959 Robert Noyce at Fairchild
Semiconductor developed a silicon integrated circuit using the planar
process, which made mass production practical. In 1965 Gordon Moore
observed that the number of components on an integrated circuit was
doubling at a regular rate, an observation that became known as Moore's law
and guided the semiconductor industry for decades.

\section{Why the stored program mattered}

The stored-program concept turned the computer from a configurable
calculator into a universal machine in Turing's sense. Because programs were
data, they could be written, copied, loaded and translated by other
programs. Assemblers, compilers and operating systems all depend on this
property. The history of software, from the subroutine library of EDSAC to
modern programming languages, begins with the decision to keep instructions
in memory.

\end{document}
